%%%%%%%%%%  scan idx 462  |  folio 449  %%%%%%%%%%
Or \(N'\) est de
rang \(d-2\), donc \(\lambda^{d-1}(N')=0\), d'où le résultat.

\subsection*{4.2.}

Soit \(i:Y\to X\) une immersion régulière. Soit \(r\) un entier \(\ge1\), et \(s:X\longrightarrow \mathbf P_X^r\) la section définie par l'homomorphisme : \((\mathcal O_X)^{r+1}\to\mathcal O_X\), passage au quotient par le sous-faisceau engendré par les \(r\) dernières coordonnées. Si \(\mathbf P_X^r\) est défini par la \(\mathcal O_X\)-Algèbre graduée \(\mathcal O_X[T_0,\ldots,T_r]\), l'idéal \(I\) définissant la section \(s\) est l'idéal engendré par les sections \(T_1,\ldots,T_r\), donc \(I/I^2\) est isomorphe à \((\mathcal O_X[T_0])^r\), donc l'idéal correspondant de \(\cO_P\) à \((\mathcal O_X)^r\).

Comme la section \(s\) est évidemment une immersion régulière, il en est de même pour \(s\circ i\), et on peut écrire la suite exacte des faisceaux conormaux correspondante (1.7). On obtient ainsi l'égalité suivante dans \(K^\circ(Y)\), en notant \(N\) le faisceau conormal de \(i\), et \(N_1\) le faisceau conormal de \(s\circ i\) :
\[
  N_1=N+r.
\]
On voit donc que pour toute immersion régulière \(i\), on peut, en la composant avec une section : \(X\to\mathbf P_X^2\), ou bien : \(X\to\mathbf P_{\mathbf P_X^1}^{1}\), obtenir une immersion régulière satisfaisant la propriété du lemme 4.1.

\textbf{Théorème 4.3.} Soient \(i:Y\to X\) une immersion fermée régulière, \(N\) la classe dans \(K^\circ(Y)\) du faisceau conormal de \(Y\) dans \(X\). On suppose que \(X\) est quasi-compact et a un faisceau ample. Alors, avec les notations de 3.3, on a pour tout \(y\in K^\circ(Y)\), et tout \(p\ge1\) :
\begin{equation}\tag{4.3.1}
 i_*\bigl(\lambda^p(N,y)\bigr)\equiv \lambda^p(i_*(y))
 \qquad\text{mod. torsion},
\end{equation}
\(\lambda^p(N,y)\) étant l'élément défini dans (V 5.3).

i) On suppose d'abord que l'hypothèse de 4.1 est satisfaite ; on va dans ce cas montrer que \(i_*\bigl(\lambda^p(N,y)\bigr)=\lambda^p(i_*(y))\). Comme \(f_*\circ f^*=\Id\), cette
